% ============================================================
%  示例：软件公司组织架构图
%  基于 harryopo-tikz-diagram org-tree 模板 v3
%  Nature 无衬线风格 — 绝对数值坐标防重叠
%
%  编译命令：xelatex example-org-tree.tex
% ============================================================

\documentclass[tikz,border=20pt]{standalone}

\usepackage{ctex}
\usepackage{tikz}
\usepackage{fontspec}
\usepackage{fontawesome5}
\usepackage{xcolor}
\usetikzlibrary{positioning, calc}

% ============================================================
% 字体设置
% ============================================================
\setmainfont{Arial}
\setsansfont{Arial}
\setmonofont{Consolas}[Scale=0.90]
\setCJKsansfont{Microsoft YaHei}
\setCJKmainfont{Microsoft YaHei}

\newcommand{\icon}[1]{{\footnotesize\color{PaperBorder}#1}}

% ============================================================
% Nature 风格配色
% ============================================================
\definecolor{PaperBorder}{HTML}{B0C4DE}
\definecolor{PaperFill}{HTML}{FFFFFF}
\definecolor{PaperTitle}{HTML}{2C3E50}
\definecolor{PaperMuted}{HTML}{5D6D7E}
\definecolor{AccentOrange}{HTML}{D35400}
\definecolor{PaperLayerBg}{HTML}{F5F8FA}

% ============================================================
%  尺寸常量（v3：精确计算防重叠）
% ============================================================
\def\DeptStep{5.0cm}    % 部门中心距（≥ 技术部团队半宽 + 产品部团队半宽）
\def\TeamW{2.2cm}       % 团队节点宽度
\def\TeamStep{2.6cm}    % 团队中心距 = TeamW + 0.4cm gap

% 部门 X 坐标（直接数值，避免 TikZ calc 问题）
\def\Xtech{-7.5}   % -1.5 * DeptStep
\def\Xprod{-2.5}   % -0.5 * DeptStep
\def\Xmkt{2.5}     %  0.5 * DeptStep
\def\Xhr{7.5}      %  1.5 * DeptStep

% 层级 Y 坐标
\def\Ytitle{0}
\def\Ytop{-1.5}
\def\Ydept{-4.0}
\def\Yteam{-6.5}
\def\Yteam2{-7.8}  % 技术部2x2第二行

\tikzset{
  org-top/.style = {
    rounded corners = 8pt,
    draw = PaperBorder,
    fill = AccentOrange!10,
    line width = 1.2pt,
    font = \normalsize\sffamily\bfseries,
    minimum width = 3.8cm,
    minimum height = 1.2cm,
    inner sep = 6pt,
    align = center,
    text = PaperTitle,
  },
  org-dept/.style = {
    rounded corners = 6pt,
    draw = PaperBorder,
    fill = PaperFill,
    line width = 0.9pt,
    font = \small\sffamily\bfseries,
    minimum width = 3.2cm,
    minimum height = 1.0cm,
    inner sep = 5pt,
    align = center,
    text = PaperTitle,
  },
  org-team/.style = {
    rounded corners = 5pt,
    draw = PaperBorder,
    fill = PaperLayerBg,
    line width = 0.7pt,
    font = \footnotesize\sffamily,
    minimum width = \TeamW,
    minimum height = 0.8cm,
    inner sep = 4pt,
    align = center,
    text = PaperTitle,
  },
  org-line/.style = {
    line width = 1.0pt,
    draw = PaperBorder,
  },
  org-title/.style = {
    font = \Large\bfseries\sffamily\color{PaperTitle},
  },
}

\begin{document}

\begin{tikzpicture}

  % ============================================================
  %  节点（全部用直接数值坐标，避免 \def 展开问题）
  % ============================================================
  \node[org-title] at (0, 0) {\icon{\faSitemap}\ 软件公司组织架构};

  % ---- 顶层：总经理 ----
  \node[org-top] (gm) at (0, -1.5) {\icon{\faUserTie}\ 总经理};

  % ---- 第二层：4个部门 ----
  \node[org-dept] (tech) at (-7.5, -4.0) {\icon{\faCode}\ 技术部};
  \node[org-dept] (prod) at (-2.5, -4.0) {\icon{\faLightbulb}\ 产品部};
  \node[org-dept] (mkt)  at ( 2.5, -4.0) {\icon{\faBullhorn}\ 市场部};
  \node[org-dept] (hr)   at ( 7.5, -4.0) {\icon{\faUsers}\ 人事部};

  % ---- 第三层：团队 ----
  % 技术部4个团队：2x2排列
  \node[org-team] (fe)  at (-8.8, -6.5) {前端};
  \node[org-team] (be)  at (-6.2, -6.5) {后端};
  \node[org-team] (qa)  at (-8.8, -7.8) {测试};
  \node[org-team] (ops) at (-6.2, -7.8) {运维};

  % 产品部2个团队
  \node[org-team] (pm)  at (-3.8, -6.5) {产品经理};
  \node[org-team] (ui)  at (-1.2, -6.5) {UI设计};

  % 市场部2个团队
  \node[org-team] (promo) at (1.2, -6.5) {推广};
  \node[org-team] (sales) at (3.8, -6.5) {销售};

  % ============================================================
  %  连接线
  % ============================================================
  % 总经理 → 部门
  \coordinate (mid1) at (0, -2.5);
  \draw[org-line] (gm.south) -- (mid1);
  \draw[org-line] (mid1) -| (tech.north);
  \draw[org-line] (mid1) -| (prod.north);
  \draw[org-line] (mid1) -| (mkt.north);
  \draw[org-line] (mid1) -| (hr.north);

  % 技术部 → 4个子团队（2x2）
  \coordinate (mid2a) at (-7.5, -5.3);
  \draw[org-line] (tech.south) -- (mid2a);
  \draw[org-line] (mid2a) -| (fe.north);
  \draw[org-line] (mid2a) -| (be.north);
  \coordinate (mid2b) at (-7.5, -7.1);
  \draw[org-line] (mid2a) -- (mid2b);
  \draw[org-line] (mid2b) -| (qa.north);
  \draw[org-line] (mid2b) -| (ops.north);

  % 产品部 → 2个子团队
  \coordinate (mid3) at (-2.5, -5.3);
  \draw[org-line] (prod.south) -- (mid3);
  \draw[org-line] (mid3) -| (pm.north);
  \draw[org-line] (mid3) -| (ui.north);

  % 市场部 → 2个子团队
  \coordinate (mid4) at (2.5, -5.3);
  \draw[org-line] (mkt.south) -- (mid4);
  \draw[org-line] (mid4) -| (promo.north);
  \draw[org-line] (mid4) -| (sales.north);

\end{tikzpicture}

\end{document}
